\section{서\ 론}
\label{sec:intro}

정간보(井間譜)는 15세기 세종대에 창안된 격자형 악보로, 우물 정(井) 자 모양의
칸(정간) 하나가 한 박을 나타내고 칸 안에 율명(律名)을 적어 음높이를 나타낸다.
음의 길이를 별도의 기호가 아니라 \emph{칸과 칸 안에서의 위치}로 표현한다는 점에서
동아시아 최초의 유량악보(mensural notation)로 평가되며~\cite{kwon2000notation,han2024sixdragons},
창안 이후 600년 가까이 지난 오늘날에도 국악 교육과정과 전문 연주 현장에서
표준 기보 체계로 쓰이고 있다.

그러나 정간보의 디지털 제작 환경은 그 위상에 크게 못 미친다.
Finale, Sibelius, MuseScore~\cite{musescore} 등 널리 쓰이는 사보 소프트웨어는
서양 오선보를 전제로 설계되어 정간보의 핵심 구조를 표현할 수 없다.
구체적으로 (1) 시가가 음표 기호가 아닌 격자 내 위치로 결정되는 점,
(2) 각(刻, 세로 열)이 오른쪽에서 왼쪽으로 진행하는 우종서(右縱書) 배치,
(3) 율명 한자·한글 표기와 옥타브 변형자(예: 청황 潢),
(4) 음표 곁에 붙는 시김새(농음·퇴성·추성 등) 장식 기호가 모두 오선보 모델
바깥에 있다\footnote{MuseScore 커뮤니티에도 정간보 지원 요청이 있으나
구현되지 않았다~\cite{musescore-jgb-request}.}.
그 결과 국악 현장에서는 한글 워드프로세서의 표 기능으로 정간보를 수작업
조판하는 방식이 지금도 일반적이다. 이렇게 만들어진 악보는 시각적 산출물일 뿐
음악 데이터가 아니므로 검색·재생·이조(移調)·재활용이 불가능하고,
조판 노동이 반복되며, 기관·개인마다 서식이 달라 교환성도 낮다.

정간보 전산화 연구가 없었던 것은 아니다. XML 문서형 정의(DTD)로 정간보를
기술하는 방법이 제안되었고~\cite{lee2013xml}, 이를 발전시킨 기술·복원 방법이
특허로 등록되었으며~\cite{patent2017jgb}, 웹 기반 정간보 프로그램의 설계
연구~\cite{jkca2022web}, 최근에는 광학악보인식(OMR)으로 정간보 인쇄 악보를
데이터화하고~\cite{kim2025omr} 트랜스포머로 15세기 궁중음악을 학습·생성하는
연구~\cite{han2024sixdragons}까지 이어졌다. 그러나 이들 연구는 인코딩 체계의
제안이나 인식·생성 모델에 초점이 있고, 국악인·교육자가 실제로 악보를
\emph{작성--조판--출판--재생}하는 전 과정을 지원하는 도구는 여전히 드물다.
공개된 웹 편집기~\cite{jeongganbo-editor-web}가 있으나 시김새·장단·가사·
조판 세부 조정 등 실무에 필요한 표현 범위에는 이르지 못했다.

본 논문은 이 간극을 메우는 웹 기반 정간보 편집·조판 시스템
``정간보 생성기''의 설계와 구현을 보고한다. 본 시스템의 기여는 다음과 같다.

\begin{itemize}
  \item \textbf{사람이 읽는 텍스트 문법}: 줄바꿈 = 각, \code{|} = 정간,
        공백 = 분박(分拍)이라는 최소 규칙으로 정간보의 격자 구조를
        일대일로 반영하는 경량 표기 문법을 설계했다(3장).
        같은 내용을 XML로 기술할 때보다 현저히 간결하면서도,
        시김새의 크기·위치 미세 조정값까지 텍스트로 직렬화되어
        버전 관리와 문서 교환에 유리하다.
  \item \textbf{이중 입력 구조}: 문법을 모르는 사용자를 위한
        클릭--입력 카드 방식(WYSIWYG)과, 대량 수정·복사에 강한
        텍스트 에디터 방식을 제공하고 두 뷰를 양방향 동기화했다(4.3절).
  \item \textbf{전통 조판 관행의 규칙화}: 우종서 각 배치, 대강(大綱)선,
        이음표(\code{-}) 행의 세로 비중 축약, 옥타브 변형 한자 우선 표기 등
        숙련 필사자의 관행을 렌더링 규칙으로 명문화하고,
        시김새 59종을 벡터(SVG)로 그린다(4.4--4.5절).
  \item \textbf{서버리스 단일 정적 앱}: 시스템 전체가 정적 파일 세 개와
        내장 데이터 파일로 구성되어 설치·계정·네트워크 없이 동작한다.
        A4 물리 단위(mm) 기반 SVG 조판에서 인쇄·PDF·300\,dpi PNG를 얻고,
        JSON 교환 포맷(\code{.jgb.json})으로 문서를 보존·공유하며,
        Web Audio로 악보를 소리 내어 확인할 수 있다(4.7--4.8절).
\end{itemize}

이하 2장에서 관련 연구를 검토하고, 3장에서 텍스트 문법을, 4장에서 시스템
설계와 구현을 기술한다. 5장에서 기존 도구·연구와 비교하고 한계를 논의한 뒤
6장에서 맺는다.
